\subsection{The complete prime-to-$p$ index}
The first theorem determines the character index at every absolute ramification index prime to $p$, in the full height-two filtration class defined above.

\begin{theorem}[The complete prime-to-$p$ character index]\label{thm:complete-prime-index}
Let $p\ge5$ and let $e$ be a positive multiple of $p-1$ with $p\nmid e$. Consider complete discretely valued characteristic-zero fields with perfect residue field of characteristic $p$, absolute index $e$, and $\mu_p\subset K$. In the height-two Barsotti--Tate filtration class,
\begin{equation}\label{eq:complete-prime-index}
 \nuMin(p,e,2)=\left\lceil\frac{1+\log_p e}{2}\right\rceil.
\end{equation}
Equivalently, $\nuMin(p,p-1,2)=1$, and for every integer $r\ge2$,
\[
 \nuMin(p,e,2)=r\qquad\text{if}\qquad p^{2r-3}<e<p^{2r-1}.
\]
Put $r=\nuMin(p,e,2)$. The upper statement is uniform over every field in the specified base class: $p^{r+1}$-torsion determines the special prime-to-$p$ inertia character. At every such $e$ there are actual allowed groups over some field with residue $\overline{\F}_p$, isomorphic $p^r$-torsion and different special characters. These lower examples are Barsotti--Tate groups; no assertion that every pair comes from elliptic curves is made.
\end{theorem}

The strict endpoints cause no ambiguity because $p\nmid e$. In particular, the complete intervals for indices two, three and four are
\[
 p<e<p^3,\qquad p^3<e<p^5,\qquad p^5<e<p^7,
\]
respectively, with the same divisibility and base-field restrictions. The proof combines an upper comparison theorem with actual lower constructions at every index.

\begin{theorem}[Iterated supersingular recovery]\label{thm:ss-iterated-range}
Let $p\ge5$, $n\ge3$, and let $K$ be a complete discretely valued field of characteristic zero with perfect residue field of characteristic $p$, containing $\mu_p$. Put $e=e_K$ and assume $p\nmid e$. If
\begin{equation}\label{eq:ss-iterated-range}
 e<\frac{(p-1)p^{3n-3}}{p^n-1},
\end{equation}
then any isomorphism $G[p^n]\simeq H[p^n]$ between height-two Barsotti--Tate groups with potentially supersingular reduction implies $\theta_G=\theta_H$. Since $(p-1)\mid e$, the numerical condition is equivalently
\[
 e\le(p-1)\bigl(p^{2n-3}+p^{n-3}\bigr).
\]
\end{theorem}

\begin{corollary}[The full filtration range]\label{cor:complete-prime-recovery}
Under the same base hypotheses, for every $n\ge3$ one has
\[
 e<p^{2n-3}\quad\Longrightarrow\quad\nuMin(p,e,2)\le n-1.
\]
\end{corollary}

Section~\ref{sec:canonical-iteration} proves these statements by repeatedly taking a canonical quotient of one model and the quotient by its transported kernel in the other. Each step loses exactly one torsion level. The graph valuation bound ensures that the first model enters the central range while at least $p^2$-torsion remains, so Theorem~\ref{thm:central-p2-rigidity} applies. The transported kernel need not be canonical in the partner. The sharp mixed obstruction then supplies both the full-class upper threshold and the matching lower examples.
